% SPDX-License-Identifier: LGPL-3.0-only
\documentclass[11pt,a4paper]{unbloated-book}

\title{Notes on Numerical Integration}
\author{Unbloated Contributors}
\date{}
\hypersetup{
  pdftitle={Notes on Numerical Integration},
  pdfauthor={Unbloated Contributors},
  pdflang={en}
}

\begin{document}
\frontmatter
\maketitle
\tableofcontents

\chapter{Preface}
This short book introduces the midpoint rule and its implementation.
The examples use semantic headings, numbered equations, and references
that remain valid when chapters move.

\mainmatter
\input{example/chapters/integration}
\input{example/chapters/implementation}

\appendix
\chapter{Notation}
\label{app:notation}
The interval endpoints are \(a\) and \(b\). Split this interval into
\(n\) equal parts with step size \(h=(b-a)/n\).

\backmatter
\begin{thebibliography}{9}
\bibitem{calculus}
OpenStax.
\emph{Calculus Volume 2}.
Rice University, 2016.
\href{https://openstax.org/details/books/calculus-volume-2}{Publisher's edition}.
\end{thebibliography}
\end{document}
